\begin{frame}{전체 그림: 4주 파이프라인}
  \begin{columns}[T]
    \begin{column}{0.6\textwidth}
      \centering
      \includegraphics[width=\linewidth]{ch08_1_rl_pipeline.png}
      \vskip0.3em
      {\footnotesize 환경·문제 정의(1주) $\to$ 알고리즘 $\geq$2개 구현·
      시드 $\geq$3·학습 곡선(2$\sim$3주) $\to$ 수식적 정당화·보고서
      (3$\sim$4주) $\to$ 발표·peer-review(4주).}
    \end{column}
    \begin{column}{0.38\textwidth}
      \begin{itemize}
        \item 흐름: \textbf{환경 $\to$ 구현·비교 $\to$ 평가·정당화 $\to$
          발표}.
        \item \kq{붉은 점선(평가 $\to$ 구현·비교)이 핵심} --- 시드 간
          편차가 비정상적으로 크거나 평가 프로토콜 결함이 발견되면
          \textbf{돌아가서 다시}.
        \item ``한 번 지나가면 돌아올 수 없는'' 것은 \textbf{최종 발표
         뿐} --- 그 외 모든 단계는 되돌아갈 수 있다.
      \end{itemize}
    \end{column}
  \end{columns}
\end{frame}
